\documentclass[10pt]{article}
\usepackage[margin=0.9in]{geometry}
\usepackage[T1]{fontenc}
\usepackage{lmodern}
\usepackage{amsmath,amssymb,amsthm}
\usepackage[hidelinks]{hyperref}
\newtheorem{theorem}{Theorem}
\newcommand{\C}{\mathbb C}
\newcommand{\D}{\mathbb D}
\newcommand{\ind}{\mathbf 1}
\newcommand{\dd}{\,dA}
\newcommand{\pv}{\operatorname{pv}}
\newcommand{\dbar}{\bar\partial}
\title{Explicit Beurling-transform counterexamples in an analytic annihilator}
\author{Alec Kriebel\thanks{Independent researcher. ORCID: \href{https://orcid.org/0009-0001-9320-500X}{0009-0001-9320-500X}.}}
\date{17 October 2026}
\begin{document}
\maketitle
\begin{abstract}
We give a self-contained expository verification of classical Beurling-transform
formulas in the analytic annihilator of Hayman--Lingham Problem~7.72.
Iwaniec's bounded angular example has an essentially unbounded logarithmic
transform. A polynomial specialization of Gill's angular-mode formula has
input norm one and transform norm $H_m=\sum_{k=1}^m k^{-1}$.
The verification includes the full analytic $L^1$ test class, principal values,
boundary traces, and essential suprema. No historical priority is claimed.
\end{abstract}

\section{Conventions, attribution, and statement}
Let $\D=\{z\in\C:|z|<1\}$ and let $dA=dx\,dy$ be unnormalized area measure.
Write $A^1(\D)$ for the holomorphic functions in $L^1(\D)$ and set
\begin{equation}\label{eq:N}
 N=\left\{v\in L^\infty(\D):
       \int_{\D}v(z)\phi(z)\dd(z)=0\quad\text{for every }\phi\in A^1(\D)\right\}.
\end{equation}
The pairing is bilinear. Extend inputs by zero outside $\D$ and use
\begin{equation}\label{eq:B}
 Bv(z)=-\frac1\pi\lim_{\varepsilon\downarrow0}
       \int_{\D\setminus B(z,\varepsilon)}
          \frac{v(\zeta)}{(z-\zeta)^2}\dd(\zeta).
\end{equation}
These are the conventions of Problem~7.72, attributed there to A.~Hinkkanen,
in the 2018 manuscript of Hayman and Lingham \cite[p.~183]{HL18}.
It asks whether $Bv$ is essentially bounded for every $v\in N$, whether
$\|Bv\|_\infty\le C\|v\|_\infty$ holds with a common constant, and whether
the latter estimate holds at least on $N\cap\C[z,\bar z]$.

The ingredients below are classical. Iwaniec \cite[p.~388, (6)--(7)]{I86}
prints the exact example $\rho=(z/\bar z)\ind_{\D}$, its transform
$(1+\log|z|^2)\ind_{\D}$, and the primitive
$(z\log|z|^2)\ind_{\D}$, with the negative kernel in \eqref{eq:B}.
Gill \cite[p.~354, Theorem~1]{G10} gives, in particular, for real
bounded $a$ supported in $(0,1)$,
\begin{equation}\label{eq:Gill}
 B\left[\frac z{\bar z}a(|z|^2)\right](z)
     =a(t)-\int_t^\infty\frac{a(s)}s\,ds,\qquad t=|z|^2,
\end{equation}
and credits Iwaniec for the logarithmic example
\cite[pp.~355--356]{G10}. The polynomial calculation here is an explicit
specialization of this formula, rather than a new transform mechanism.
The 2019 book's Update~7.72 credits Volberg with resolving the problem and
refers to Ba\~nuelos--Janakiraman \cite{BJ08} for an account
\cite[p.~211]{HL19}. We report that published credit; our reading of
the cited finite-$p$ paper did not locate an explicit account of both
endpoint assertions, and we do not identify a separate Volberg proof.
The purpose here is to verify the explicit
consequences under \eqref{eq:N}--\eqref{eq:B}, independently of that history.

\begin{theorem}\label{thm:main}
Define $\rho(z)=z/\bar z$ for $0<|z|<1$, with $\rho(0)=0$.
Then $\rho\in N$, $\|\rho\|_{L^\infty(\D)}=1$, and
\begin{equation}\label{eq:rho}
 B\rho(z)=1+\log|z|^2\qquad(0<|z|<1).
\end{equation}
In particular, $B\rho\notin L^\infty(\D)$.
For every positive integer $m$, let
\begin{equation}\label{eq:vm}
 P_m(t)=\sum_{j=0}^{m-1}(1-t)^j,\qquad
 v_m(z)=z^2P_m(|z|^2),\qquad H_m=\sum_{k=1}^m\frac1k.
\end{equation}
Then $v_m\in N\cap\C[z,\bar z]$ and
\begin{equation}\label{eq:norms}
 \|v_m\|_{L^\infty(\D)}=1,\qquad
 \|Bv_m\|_{L^\infty(\D)}=H_m.
\end{equation}
Both transforms vanish almost everywhere outside $\D$.
\end{theorem}

\section{Angular cancellation and the Cauchy potential}
For bounded measurable $a:(0,1)\to\C$, put
$f_a(re^{i\theta})=e^{2i\theta}a(r^2)$ in $\D$ and zero outside.
If $\phi\in A^1(\D)$, then $f_a\phi$ is absolutely integrable.
On each circle $r<1$, the Taylor series $\phi(z)=\sum_{n\ge0}c_nz^n$
converges uniformly, and therefore
\[
 \int_0^{2\pi}e^{2i\theta}\phi(re^{i\theta})\,d\theta
  =\sum_{n\ge0}c_nr^n\int_0^{2\pi}e^{i(n+2)\theta}\,d\theta=0.
\]
Polar-coordinate Fubini proves $f_a|_{\D}\in N$. This verifies every
analytic $L^1$ test, without a density assertion about endpoint dual spaces.

For bounded compactly supported $f$, define the actual Cauchy potential
\begin{equation}\label{eq:C}
 Cf(z)=\frac1\pi\int_{\C}\frac{f(\zeta)}{z-\zeta}\dd(\zeta).
\end{equation}
This integral is absolutely convergent for every $z$, since
$|z-\zeta|^{-1}$ is locally integrable. Fix $z\ne0$ and $r\ne|z|$.
For $r<|z|$, the geometric expansion of $(z-re^{i\theta})^{-1}$ has
terms $r^ne^{in\theta}/z^{n+1}$, so multiplication by $e^{2i\theta}$
leaves no constant angular term. For $r>|z|$, use
\[
 \frac1{z-re^{i\theta}}
   =-\sum_{n\ge0}\frac{z^n}{r^{n+1}}e^{-i(n+1)\theta}.
\]
Only $n=1$ survives; its angular integral is $-2\pi z/r^2$.
Each series converges uniformly on its fixed circle. Absolute
integrability of \eqref{eq:C} permits polar Fubini and omission of the
single radius $r=|z|$. Substitution $s=r^2$ gives
\begin{equation}\label{eq:Ca}
 Cf_a(z)=
 \begin{cases}
 -z\displaystyle\int_{|z|^2}^1\frac{a(s)}s\,ds,&0<|z|<1,\\[1ex]
 0,&|z|>1.
 \end{cases}
\end{equation}
At $z=0$ the angular integral is zero. The same calculation gives zero
on $|z|=1$. Thus no unspecified holomorphic correction occurs in this
potential.

We record why differentiating it yields the principal value in
\eqref{eq:B}. With
$\partial=\tfrac12(\partial_x-i\partial_y)$ and
$\dbar=\tfrac12(\partial_x+i\partial_y)$, the locally integrable
kernel $E(z)=1/(\pi z)$ satisfies, in distributions,
\begin{equation}\label{eq:kernel}
 \dbar E=\delta_0,\qquad
 \partial E=-\frac1\pi\pv\frac1{z^2}.
\end{equation}
Indeed, integrate by parts off $|z|\le\varepsilon$ against a smooth
compactly supported test function $\psi$. For $\partial E$, the inner
boundary term has angular factor $e^{-2i\theta}$ and tends to zero.
For $\dbar E$ the factor is one and the normalized term tends to
$\psi(0)$. The principal-value distribution exists because its constant
angular term cancels and $\psi(z)-\psi(0)=O(|z|)$ is integrable against
$|z|^{-2}dA$ near zero. This establishes both the sign and normalization.

Convolution with compactly supported $f$ now gives
$\partial Cf=Bf$ in distributions. Where $f$ is locally $C^1$, the
pointwise principal value exists: on a small centered disc subtract
$f(z)$, whose kernel integral is zero, and use
\[
 \frac{|f(\zeta)-f(z)|}{|z-\zeta|^2}\le\frac{C}{|z-\zeta|}.
\]
On smaller compact neighborhoods these local integrals converge uniformly;
the remaining integral is continuous. Thus the principal-value function
is continuous there and agrees with the distributional derivative.
This applies everywhere inside $\D$ to $v_m$ and everywhere in
$\D\setminus\{0\}$ to $\rho$.

\section{The logarithmic and polynomial calculations}
Taking $a=1$ in \eqref{eq:Ca} gives
\[
 C\rho(z)=u(z)=z\log|z|^2\quad(0<|z|<1),
 \qquad u(0)=0,\qquad u(z)=0\quad(|z|\ge1).
\]
This is continuous at the origin and has zero boundary trace on
$|z|=1$. Direct differentiation in the punctured disc gives
\begin{equation}\label{eq:uderivatives}
 \dbar u=\frac z{\bar z},\qquad \partial u=1+\log|z|^2.
\end{equation}
These are also the global weak derivatives, extended by zero outside
$\D$. To check this, integrate by parts on
$\varepsilon<|z|<1$. The outer boundary term is zero. The inner term
is bounded by a constant times
$\varepsilon\sup_{|z|=\varepsilon}|u(z)|
 =2\varepsilon^2|\log\varepsilon|\to0$.
The displayed derivatives belong to every finite $L^p$ because
$\int_0^1 r|\log r|^p\,dr<\infty$; hence passage to the limit is
legitimate and creates neither a point mass nor a boundary measure.
Equations \eqref{eq:kernel} and \eqref{eq:uderivatives}, with the local
principal-value argument, prove \eqref{eq:rho}.

For $M\ge1$, its level set has positive area:
\begin{equation}\label{eq:level}
 \bigl|\{z\in\D:|B\rho(z)|>M\}\bigr|=\pi e^{-M-1}>0.
\end{equation}
Consequently the unboundedness is essential. At the exceptional point
$0$, the truncated integral is
$-\pi^{-1}\int_{\varepsilon<|\zeta|<1}|\zeta|^{-2}\dd(\zeta)
 =2\log\varepsilon\to-\infty$; no finite principal value there is
needed for the almost-everywhere conclusion. Angular cancellation above
places $\rho$ in $N$, and $|\rho|=1$ almost everywhere in $\D$.

For the polynomials in \eqref{eq:vm}, the radial amplitude is
\[
 a_m(t)=tP_m(t)=1-(1-t)^m\qquad(0<t<1).
\]
It lies between zero and one and tends to one as $t\uparrow1$.
Every level below one is exceeded on a positive-area outer annulus,
so $\|v_m\|_\infty=1$. The finite expression $z^2P_m(z\bar z)$
is a polynomial at the origin as well, and angular cancellation gives
$v_m\in N$.
Formula \eqref{eq:Ca} becomes $Cv_m(z)=zF_m(|z|^2)$ in $\D$, where
\begin{equation}\label{eq:F}
 F_m(t)=-\int_t^1P_m(s)\,ds
       =-\sum_{k=1}^m\frac{(1-t)^k}{k}.
\end{equation}
Its zero extension is continuous across $|z|=1$ and smooth at the
origin. Integration by parts, using the zero boundary value, gives its
global weak derivatives without boundary measures. Inside $\D$ they are
\[
 \dbar Cv_m=z^2P_m(|z|^2)=v_m,\qquad
 Bv_m(z)=\partial Cv_m(z)=G_m(|z|^2),
\]
where
\begin{equation}\label{eq:G}
 G_m(t)=1-(1-t)^m-\sum_{k=1}^m\frac{(1-t)^k}{k}.
\end{equation}
The output is continuous in $\D$. Its endpoint values and derivative are
\[
 G_m(0)=-H_m,\qquad G_m(1)=1,\qquad
 G_m'(t)=P_m(t)+m(1-t)^{m-1}>0\quad(0<t<1).
\]
Thus $|G_m(t)|\le H_m$, since $H_m\ge1$. Conversely, continuity at
zero makes $|G_m(|z|^2)|>H_m-\eta$ on a disc of positive area for
every $\eta>0$ with $H_m-\eta>0$. This proves the essential norm
identity in \eqref{eq:norms}, rather than relying on an isolated point
value. For $m=1$ one obtains $v_1=z^2$ and $G_1(t)=2t-1$, so the
boundary case has norm one as asserted. Since
$H_m\ge\log(m+1)\to\infty$, no common polynomial bound exists.
Finally, the potentials vanish on the exterior, so their derivatives,
and hence both transforms, vanish there. This completes the proof of
Theorem~\ref{thm:main}.

Each fixed polynomial in this family has a bounded transform; the
polynomial conclusion concerns the failure of one constant for all
inputs. The separate classical input $\rho$ disproves boundedness of
the transform for every member of $N$. The same conclusions hold if
the output norm is taken on $\C$.

\paragraph{Verification and disclosure.}
The all-parameter analytic proof above supplies the conclusions.
Finite exact-coefficient checks in the accompanying verification
material are diagnostic supplements, not a proof for all $m$ or a
formal proof-assistant certification. AI tools were used extensively
in research, solving, verification, literature investigation, drafting,
and internal adversarial review. This note is unrefereed and has had
no external human peer review.

\begin{thebibliography}{9}
\bibitem{BJ08}
R.~Ba\~nuelos and P.~Janakiraman,
\emph{$L^p$-bounds for the Beurling--Ahlfors transform},
Trans. Amer. Math. Soc. \textbf{360} (2008), no.~7, 3603--3612.
\href{https://doi.org/10.1090/S0002-9947-08-04537-6}{doi:10.1090/S0002-9947-08-04537-6}.

\bibitem{G10}
J.~T.~Gill,
\emph{On the Beurling--Ahlfors transform's weak-type constant},
Michigan Math. J. \textbf{59} (2010), no.~2, 353--363.
\href{https://doi.org/10.1307/mmj/1281531461}{doi:10.1307/mmj/1281531461}.

\bibitem{HL18}
W.~K.~Hayman and E.~F.~Lingham,
\emph{Research Problems in Function Theory (New Edition)},
2018, arXiv:1809.07200v2, Problem~7.72, printed p.~183.
\href{https://arxiv.org/abs/1809.07200v2}{arxiv.org/abs/1809.07200v2}.

\bibitem{HL19}
W.~K.~Hayman and E.~F.~Lingham,
\emph{Research Problems in Function Theory: Fiftieth Anniversary Edition},
Springer, Cham, 2019, Update~7.72, p.~211.
\href{https://doi.org/10.1007/978-3-030-25165-9}{doi:10.1007/978-3-030-25165-9}.

\bibitem{I86}
T.~Iwaniec,
\emph{The best constant in a BMO-inequality for the Beurling--Ahlfors transform},
Michigan Math. J. \textbf{33} (1986), no.~3, 387--394.
\href{https://doi.org/10.1307/mmj/1029003418}{doi:10.1307/mmj/1029003418}.
\end{thebibliography}
\end{document}
